\section{本讲的问题：GPU 为什么既快又难用？}
Lecture 5 把 Lecture 2 的 resource accounting 落到硬件：为什么 GPU 能支撑 scaling？为什么同样是 FLOPs，有的操作飞快，有的操作很慢？为什么 FlashAttention 这种算法能改变 attention 的实际速度？本讲的核心问题因此不是峰值 FLOP/s，而是算子能否把足够多的数据复用和并行工作送到执行单元，同时避免 HBM、调度与对齐成为瓶颈。
\begin{knowledgebox}{术语消化：GPU memory hierarchy}
HBM 是 High Bandwidth Memory，GPU 上的大容量高带宽显存，容量大但离计算单元远；SRAM 是 Static Random-Access Memory，常对应片上 shared memory/cache，容量小但速度快；register 是每个 thread 私有的最快存储。GPU 优化的核心之一，就是让数据从 HBM 读进来以后尽量在 SRAM/register 附近复用，而不是频繁写回 HBM。
\end{knowledgebox}
\begin{importantbox}{本讲主线}
GPU performance 的核心不是“GPU 很快”四个字，而是让数据移动、并行执行和矩阵硬件匹配。低精度、fusion、recomputation、coalescing、tiling、FlashAttention 都是在减少 HBM 往返或提高数据复用。
\end{importantbox}
